% =============================================================================
%  Chapter 17 — Polynomial Patch Recovery for TET10 Visualisation
% =============================================================================

\chapter{Polynomial Patch Recovery for TET10 Visualisation}
\label{ch:patch_recovery}

This chapter documents the second-stage fix for the TET10 visualisation path.
The first fix\,---\,avoiding unsafe lumped $L^2$ projection\,---\,was
necessary but not sufficient.  Replacing signed lumping with a plain
volume-weighted average removed the gross artefact, but it also produced a
field that was too diffusive for smooth contour rendering.  The proper
solution is a local polynomial reconstruction over nodal patches.


% ─────────────────────────────────────────────────────────────────────────
\section{Why Averaging Was Not Enough}
\label{sec:patch_not_enough}

The TET10 strain field in the cantilever benchmark is smooth at the level of
Gauss-point samples: the displacement solution is correct, the kinematics are
correct, and the constitutive update is correct.  The visual defect appears
only after transfer to the nodes.

Volume-weighted averaging fixes the sign problem but throws away too much
information:
\begin{itemize}[nosep]
  \item it compresses each element to a single average value before nodal
        transfer;
  \item it cannot reproduce linear or quadratic intra-element variation;
  \item it therefore damps gradients that should remain visible in the final
        contour plot.
\end{itemize}

For higher-order simplex elements this is the wrong trade-off.  The data are
already rich enough to support a better local fit, so a production-quality
visualisation path should use them.


% ─────────────────────────────────────────────────────────────────────────
\section{Recovered Quantity and Patch Topology}
\label{sec:patch_quantity}

Let node $I$ have coordinate $\bx_I$.  Its patch is the set of all elements
incident on that node:
\begin{equation}
  \mathcal{P}_I = \{ e \in \mathcal{E} \mid I \in e \}.
\end{equation}

The exporter gathers every Gauss point of every element in $\mathcal{P}_I$.
Each observation contributes:
\begin{equation}
  (\bx_g,\; \sigma_g,\; \omega_g),
\end{equation}
where $\bx_g$ is the physical coordinate of the sample, $\sigma_g$ is the
material field component value, and $\omega_g = |w_g J_g|$ is the positive
sample weight.

\paragraph{Important distinction.}
The recovery patch is a \emph{post-processing construct}.  It is not part of
the FE approximation space, does not add new unknowns, and does not alter the
element stiffness formulation.  It exists solely to reconstruct a nodal field
for rendering.


% ─────────────────────────────────────────────────────────────────────────
\section{Centered and Scaled Basis}
\label{sec:patch_basis}

To reduce conditioning problems and keep the constant coefficient equal to the
recovered nodal value, the polynomial basis is centered at the node:
\begin{equation}
  \hat{\bx}_g = \frac{\bx_g - \bx_I}{h_I},
  \qquad
  h_I = \max_{g \in \mathcal{P}_I} \|\bx_g - \bx_I\|_2.
\end{equation}

For the 3D case, the quadratic basis is
\begin{equation}
  \mathbf{p}(\hat{\bx}) =
  \begin{bmatrix}
  1 &
  \hat{x} &
  \hat{y} &
  \hat{z} &
  \hat{x}^2 &
  \hat{x}\hat{y} &
  \hat{x}\hat{z} &
  \hat{y}^2 &
  \hat{y}\hat{z} &
  \hat{z}^2
  \end{bmatrix}^{\!\mathsf{T}}.
\end{equation}

With this choice,
\begin{equation}
  \tilde{\sigma}_I = p_I(\bx_I) = a_{I,0},
\end{equation}
because all non-constant basis terms vanish at the origin of the local patch
coordinates.


% ─────────────────────────────────────────────────────────────────────────
\section{Least-Squares System}
\label{sec:patch_lsq}

For one field component, the recovery solves
\begin{equation}
  \ba_I
  = \arg\min_{\ba}
    \sum_{g \in \mathcal{P}_I}
    \omega_g
    \left(
      \mathbf{p}(\hat{\bx}_g)^{\mathsf{T}} \ba - \sigma_g
    \right)^2.
\end{equation}

The implementation accumulates the normal equations
\begin{equation}
  \mathbf{N}_I = \sum_g \omega_g\, \mathbf{p}_g \mathbf{p}_g^{\mathsf{T}},
  \qquad
  \mathbf{r}_I = \sum_g \omega_g\, \mathbf{p}_g \sigma_g,
\end{equation}
and solves
\begin{equation}
  \mathbf{N}_I \ba_I = \mathbf{r}_I.
\end{equation}

\paragraph{Why use \texorpdfstring{$|w_g J_g|$}{|wJ|}?}
In numerical integration the sign of $w_g$ matters because it contributes to
an integral approximation.  In patch recovery the role of the weight is
different: it measures the local influence of a sample.  A negative
quadrature weight should therefore \emph{not} negate the influence of a
perfectly valid sample in the least-squares fit.  Using $|w_g J_g|$ is the
correct semantics for recovery.


% ─────────────────────────────────────────────────────────────────────────
\section{Rank-Adaptive Fallback}
\label{sec:patch_rank}

Not every nodal patch can support a full quadratic fit:
\begin{itemize}[nosep]
  \item a boundary node may see too few neighbouring elements;
  \item a coarse mesh may provide too few independent Gauss points;
  \item a highly directional patch may lead to rank deficiency.
\end{itemize}

The exporter therefore uses the following hierarchy:
\begin{enumerate}[nosep]
  \item quadratic fit if the patch has enough samples and full rank;
  \item linear fit if the quadratic system is rank-deficient;
  \item weighted average if even the linear fit is not admissible.
\end{enumerate}

This hierarchy matters in practice.  A single TET10 contributes four HMS Gauss
points, which is not enough for a quadratic 3D fit (ten coefficients) but is
exactly enough for a linear fit (four coefficients).  The new regression tests
check this explicitly.


% ─────────────────────────────────────────────────────────────────────────
\section{Implementation Choices and HPC Considerations}
\label{sec:patch_hpc}

The implementation in \texttt{VTKModelExporter.hh} was written to avoid
contaminating the assembly kernels with post-processing concerns:
\begin{itemize}[nosep]
  \item all recovery logic is isolated in the exporter;
  \item the FE solve path, material update path, and element geometry kernels
        are unchanged;
  \item the recovery loop uses geometry already bound in the elements, so no
        new ownership or pointer chains were introduced into the hot path;
  \item the least-squares accumulation uses fixed-size small Eigen matrices for
        the normal equations, with only a tiny dynamic solve on the active
        block size.
\end{itemize}

This is the right architectural compromise for \texttt{fall\_n}:
\begin{itemize}[nosep]
  \item \textbf{compile-time and runtime discipline} are preserved in the FE
        kernels;
  \item \textbf{visualisation quality} is upgraded where it belongs, in
        post-processing;
  \item \textbf{traceability} remains high because the policy is explicit in
        \texttt{MaterialFieldProjection}.
\end{itemize}


% ─────────────────────────────────────────────────────────────────────────
\section{Regression Coverage}
\label{sec:patch_tests}

Two new tests were added in \texttt{tests/test\_simplex.cpp}:
\begin{enumerate}[nosep]
  \item \texttt{test\_tet10\_patch\_recovery\_recovers\_linear\_field\_on\_single\_element()}
        verifies that the fallback from quadratic to linear fit works and is
        exact for a linear field on one straight-sided TET10.
  \item \texttt{test\_tet10\_patch\_recovery\_recovers\_quadratic\_field\_on\_patch()}
        verifies exact recovery of a quadratic scalar field at a shared patch
        node using multiple TET10 elements and real simplex quadrature points.
\end{enumerate}

These tests are more valuable than a screenshot-based validation because they
check the transfer operator directly, independently of ParaView rendering
settings.


% ─────────────────────────────────────────────────────────────────────────
\section{Architectural Outcome}
\label{sec:patch_outcome}

The final picture is now clean:
\begin{itemize}[nosep]
  \item \textbf{HEX27} stays on lumped $L^2$ because its local nodal lumping is
        positive and well behaved.
  \item \textbf{TET10} automatically switches to polynomial patch recovery
        because signed lumping makes row-sum $L^2$ unsuitable as a nodal
        average.
  \item \textbf{Volume averaging} remains available for debugging and for the
        rare case where a deliberately diffusive fallback is preferred.
\end{itemize}

This resolves the remaining mismatch between mathematically correct Gauss-point
fields and visually acceptable nodal contours without degrading the
performance-oriented architecture of the main finite element code.
